% 第12章 音響ピークの解剖学（WP5 skeleton; 本文・導出は執筆予定）
\section{音響ピークの解剖学}
\label{ch:ch12_peaks}

\paragraph{この章で導くこと（入力→出力）}
解析理解と射影を合成し、SWプラトー・ピーク位置・奇偶非対称・減衰尾・ISW を指差し確認する。

% --- 節構成（01_textbook_spec.md より）-------------------------------------
\subsection{ピーク位置の公式（$\ell_1\approx220$）}
% TODO(WP5): 12.1 の本文。
\subsection{奇数/偶数ピークの非対称}
% TODO(WP5): 12.2 の本文。
\subsection{SWプラトーと積分SW（契約 S12.1）}
% TODO(WP5): 12.3 の本文。
\subsection{減衰尾}
% TODO(WP5): 12.4 の本文。
\subsection{総合図解}
% TODO(WP5): 12.5 の本文。


\paragraph{導出契約}
本章の導出契約: S12.1。各契約は「入力式→出力式」を中間ステップ省略なし
（1ステップ＝学部4年生が5分で検算可能）で履行する。\emph{本文プローズは WP5 で執筆。}

\paragraph{まとめ・チェックリスト・演習}
% TODO(WP5): 章末まとめ、チェックリスト、演習 3–6 問（うち1問は対応 NB の課題）。
